\scanpage{185}
\mspage{II.33}
% archive sheet number: 184
% block: top, lines 1-2 | text
Cela se ramène formellement \deleted{au} \deleted{\unsure{cas où $\pi$ est un groupe analytiquement} \illegible{2}}\inserted{\unsure{au thm bien connu}}

% block: upper third, lines 3-5 | text
\uline{\unsure{Lemme|Théorème}} Soit $G$ un groupe analytique, $\mathfrak g$ son alg. de Lie, qui est aussi l'alg. de Lie de ses sous-groupes ouverts \inserted{$U_i$}. Alors le foncteur

% block: upper third, line 6 | display
\[
\unsure{\varinjlim_i|\underset{i}{\underline{\lim}}}\operatorname{Mod}_{\mathbb Q_\ell}U_i\longrightarrow\operatorname{Mod}_{\deleted{\unsure{\mathbb Q_\ell}}}\mathfrak g
\]

% block: upper third, line 7 | text
est une équivalence de catégories.

% block: middle third, lines 8-13 | text
6.5. On fera attention \deleted{\unsure{qu'}}\inserted{que même si $\pi=\widehat{\mathbb Z}$,} on ne récupère pas $\pi$ (ou les \unsure{proobjets} $G$ \inserted{\unsure{qu'on vient associer}} (ni même le pro-Lie $\mathfrak g$) par la connaissance de la catégorie $\operatorname{Mod}_{\mathbb Q_\ell}(\pi)$, même munie de sa structure de produit tensoriel! on ne récupère que des enveloppes \uline{algébriques}. [\unsure{On montrera} cependant si \deleted{$\unsure{G}$} les $\mathfrak g_\alpha$ sont semi simples\ldots].

% block: lower left margin beside construction | margin note
% slanted
\margin{\fbox{\parbox{25mm}{Attention, il est essentiel ici de travailler sur $\mathbb Q_\ell$, non sur $\mathbb Z_\ell$; \unsure{le faire réserve des surprises}!}}}

% block: middle to lower third, lines 14-18 | text
\deleted{6.6} Je vais construire un groupe pro-\uline{algébrique} \inserted{\uline{affine}} canonique \inserted{$\pi^{\mathrm{alg}}$} sur \textcircled{$\mathbb Q_\ell$}, dont les représentations algébriques sur \deleted{$\unsure{\pi}$} $\mathbb Q_\ell$ vont définir la catégorie $\operatorname{Mod}_{\mathbb Q_\ell}(\pi)$ à équivalence près, en respectant la structure tensorielle. Pour ceci, on considère, pour tt $E\in\operatorname{Ob}\operatorname{Mod}_{\mathbb Q_\ell}(\pi)$, le groupe algébrique engendré dans $\underline{\operatorname{Aut}}_{\mathbb Q_\ell}(E)$ par l'image de $\pi$, soit $\widetilde G_E$. S'il existe un hom \uline{surjectif} $E\mathbin{-}F$, il y a un hom naturel $\widetilde G_E\mathbin{-}\widetilde G_F$\inserted{qui est un épimorphisme (\unsure{donc} surjectif)}, et de cette façon les $(G_E)$
